\uuid{U3cX}
\exo7id{3941}
\titre{exo7 3941}
\auteur{quercia}
\organisation{exo7}
\datecreate{2010-03-11}
\isIndication{false}
\isCorrection{false}
\chapitre{Dérivabilité des fonctions réelles}
\sousChapitre{Autre}
\cpge{mpsi-mp2i.raisonnement-vocabulaire-ensembliste.c, pcsi.raisonnement-vocabulaire-ensembliste.d, ptsi.raisonnement-vocabulaire-ensembliste.d, mpsi-mp2i.fonctions-valeurs-reelles-complexes.b, pcsi.techniques-calcul-fonctions.b, ptsi.techniques-calcul-fonctions.b, mpsi-mp2i.limites-continuite.c, pcsi.limites-continuite.c, ptsi.limites-continuite.c, bcpst1.fonctions-reelles-usuelles.a, bcpst1.limites-continuite.c, bcpst1.derivation.a}
\module{Analyse}
\niveau{L1}
\difficulte{}

\contenu{
\texte{
Soit $f : {[a,b]} \to \R$ dérivable.
}
\begin{enumerate}
    \item \question{On note $E = \{(x,y) \in {[a,b]}^2 \text{ tq } x < y\}$ et pour $(x,y) \in E$ :
    $\varphi(x,y) = \frac{f(x)-f(y)}{x-y}$.
    Montrer que $\varphi(E)$ est un intervalle.}
    \item \question{En déduire que $f'([a,b])$ est un intervalle.}
\end{enumerate}
}
